\documentclass[11pt,a4paper]{article}
\usepackage[spanish,es-noquoting]{babel}
\usepackage[top=2.2cm, bottom=2.2cm, left=2.2cm, right=2.2cm]{geometry}
\usepackage{fontspec}
\usepackage{amsmath}
\usepackage[version=4]{mhchem}
\usepackage{enumitem}
\usepackage{microtype}

\setmainfont{Calibri}
\linespread{1.15}

\begin{document}
\pagestyle{empty}

% PÁGINA 1: RESOLUCIÓN 1
\begin{center}
    {\LARGE \textbf{Resolución: Ejercicio 1 (Combustión de Propano)}}\\[4pt]
    {\large \textbf{Química 4.º ESO / 1.º de Bachillerato}}
\end{center}

\vspace{0.4cm}
\noindent\textbf{Cálculo de masas molares ($M$):}
\begin{itemize}[leftmargin=1.5em, itemsep=2pt]
    \item $M(\ce{C3H8}) = 3(12) + 8(1) = 44\text{ g/mol}$
    \item $M(\ce{CO2}) = 12 + 2(16) = 44\text{ g/mol}$
    \item $M(\ce{O2}) = 32\text{ g/mol} \qquad M(\ce{H2O}) = 18\text{ g/mol}$
\end{itemize}

\vspace{0.2cm}
\begin{enumerate}[label=\textbf{\alph*)}, leftmargin=2em, itemsep=10pt]
    \item \textbf{Ajuste de la reacción química:}
    \[
    \ce{C3H8 (g) + 5 O2 (g) -> 3 CO2 (g) + 4 H2O (l)}
    \]

    \item \textbf{Moles y masa teórica máxima de \ce{CO2} ($100\%$ de rendimiento):}
    \begin{itemize}[leftmargin=1.5em, itemsep=3pt]
        \item Moles de reactivo inicial (\ce{C3H8}):
        \[
        n(\ce{C3H8}) = \frac{88\text{ g}}{44\text{ g/mol}} = 2{,}00\text{ mol de \ce{C3H8}}
        \]
        \item Relación estequiométrica ($1\text{ mol \ce{C3H8}} : 3\text{ mol \ce{CO2}}$):
        \[
        n_{\text{teórico}}(\ce{CO2}) = 2{,}00\text{ mol \ce{C3H8}} \times \frac{3\text{ mol \ce{CO2}}}{1\text{ mol \ce{C3H8}}} = \mathbf{6{,}00\text{ mol de \ce{CO2}}}
        \]
        \item Masa teórica de \ce{CO2}:
        \[
        m_{\text{teórica}}(\ce{CO2}) = 6{,}00\text{ mol} \times 44\text{ g/mol} = \mathbf{264\text{ g de \ce{CO2}}}
        \]
    \end{itemize}

    \item \textbf{Masa real de \ce{CO2} emitida con un rendimiento del $85\%$:}
    \[
    m_{\text{real}}(\ce{CO2}) = m_{\text{teórica}} \times 0{,}85 = 264\text{ g} \times 0{,}85 = \mathbf{224{,}4\text{ g de \ce{CO2}}}
    \]

    \item \textbf{Volumen real de \ce{CO2} en Condiciones Normales ($\text{C.N.}$):}
    \begin{itemize}[leftmargin=1.5em, itemsep=3pt]
        \item Moles reales de \ce{CO2} producidos al $85\%$:
        \[
        n_{\text{real}}(\ce{CO2}) = 6{,}00\text{ mol} \times 0{,}85 = 5{,}10\text{ mol de \ce{CO2}}
        \]
        \item Volumen en C.N. ($V_{\text{molar}} = 22{,}4\text{ L/mol}$):
        \[
        V_{\text{real}}(\ce{CO2}) = 5{,}10\text{ mol} \times 22{,}4\text{ L/mol} = \mathbf{114{,}24\text{ L de \ce{CO2}}}
        \]
    \end{itemize}
\end{enumerate}

\newpage

% PÁGINA 2: RESOLUCIÓN 2
\begin{center}
    {\LARGE \textbf{Resolución: Ejercicio 2 (Reducción de \ce{Fe2O3} en Alto Horno)}}\\[4pt]
    {\large \textbf{Química 4.º ESO / 1.º de Bachillerato}}
\end{center}

\vspace{0.4cm}
\noindent\textbf{Cálculo de masas molares ($M$):}
\begin{itemize}[leftmargin=1.5em, itemsep=2pt]
    \item $M(\ce{Fe2O3}) = 2(56) + 3(16) = 112 + 48 = 160\text{ g/mol}$
    \item $M(\ce{Fe}) = 56\text{ g/mol} \qquad M(\ce{CO2}) = 44\text{ g/mol}$
\end{itemize}

\vspace{0.2cm}
\begin{enumerate}[label=\textbf{\alph*)}, leftmargin=2em, itemsep=10pt]
    \item \textbf{Ajuste de la ecuación química:}
    \[
    \ce{Fe2O3 (s) + 3 CO (g) -> 2 Fe (s) + 3 CO2 (g)}
    \]

    \item \textbf{Masa teórica máxima de hierro puro (\ce{Fe}) al $100\%$:}
    \begin{itemize}[leftmargin=1.5em, itemsep=3pt]
        \item Moles iniciales de óxido férrico:
        \[
        n(\ce{Fe2O3}) = \frac{320\text{ g}}{160\text{ g/mol}} = 2{,}00\text{ mol de \ce{Fe2O3}}
        \]
        \item Moles teóricos de hierro ($1\text{ mol \ce{Fe2O3}} : 2\text{ mol \ce{Fe}}$):
        \[
        n_{\text{teórico}}(\ce{Fe}) = 2{,}00\text{ mol \ce{Fe2O3}} \times \frac{2\text{ mol \ce{Fe}}}{1\text{ mol \ce{Fe2O3}}} = \mathbf{4{,}00\text{ mol de \ce{Fe}}}
        \]
        \item Masa teórica de \ce{Fe}:
        \[
        m_{\text{teórica}}(\ce{Fe}) = 4{,}00\text{ mol} \times 56\text{ g/mol} = \mathbf{224\text{ g de \ce{Fe}}}
        \]
    \end{itemize}

    \item \textbf{Porcentaje de rendimiento ($\%$) del proceso:}
    \[
    \text{Rendimiento} = \frac{m_{\text{real}}}{m_{\text{teórica}}} \times 100 = \frac{179{,}2\text{ g}}{224\text{ g}} \times 100 = \mathbf{80\%}
    \]

    \item \textbf{Masa real de dióxido de carbono (\ce{CO2}) generada:}
    \begin{itemize}[leftmargin=1.5em, itemsep=3pt]
        \item Moles teóricos de \ce{CO2} ($1\text{ mol \ce{Fe2O3}} : 3\text{ mol \ce{CO2}}$):
        \[
        n_{\text{teórico}}(\ce{CO2}) = 2{,}00\text{ mol \ce{Fe2O3}} \times \frac{3\text{ mol \ce{CO2}}}{1\text{ mol \ce{Fe2O3}}} = 6{,}00\text{ mol de \ce{CO2}}
        \]
        \item Moles reales con el $80\%$ de rendimiento:
        \[
        n_{\text{real}}(\ce{CO2}) = 6{,}00\text{ mol} \times 0{,}80 = 4{,}80\text{ mol de \ce{CO2}}
        \]
        \item Masa real generada:
        \[
        m_{\text{real}}(\ce{CO2}) = 4{,}80\text{ mol} \times 44\text{ g/mol} = \mathbf{211{,}2\text{ g de \ce{CO2}}}
        \]
    \end{itemize}
\end{enumerate}

\end{document}
